\documentclass[10pt,a4paper]{amsart}
\usepackage[margin=19mm]{geometry}
\usepackage{amsmath,amssymb,mathtools}
\usepackage[scr=rsfs,cal=euler]{mathalfa}
\usepackage{xfrac}
\usepackage[unicode,hidelinks,bookmarks=false]{hyperref}
\newcommand{\ie}{i.e.\ }
\newcommand{\cf}{cf.\ }
\newcommand{\resp}{respectively\ }
\newcommand{\Subsection}{\subsection}
\providecommand{\nobreakdash}{\nobreak\hbox{-}}
\setlength{\emergencystretch}{2em}
\multlinegap0pt
\newtheorem{proposition}{Proposition}
\newcommand{\F}{\mathbb{F}}
\newcommand{\Tr}{\mathrm{Tr}}
\newcommand{\charac}{\operatorname{char}}
\title{Additive Conjucyclic Codes over $\F_{q^2}$: A Trace Correspondence and Quantum Error-Correction}
\author{Jingjie Lv, Xian Lian, Ruihu Li, Hanxu Hou}
\date{Source: arXiv:2007.03963v3 (CC BY 4.0)}
\begin{document}
\maketitle

\noindent\emph{Source and licence.} Additive Conjucyclic Codes over $\F_{q^2}$: A Trace Correspondence and Quantum Error-Correction. Version arXiv:2007.03963v3 (CC BY 4.0), distributed by arXiv under the Creative Commons Attribution 4.0 International licence, http://creativecommons.org/licenses/by/4.0/, as recorded in the arXiv licence metadata for this exact version and shown on the abstract page https://arxiv.org/abs/2007.03963v3. Full licence notice: \texttt{licenses/arxiv-2007.03963-CC-BY-4.0.txt}.\par\smallskip

\noindent\textit{Editorial scope.} Counted unit: the article's proposition proposition (source0.tex lines 209--211). The article's complete original proof of it is delivered with it (source0.tex lines 212--257, one \texttt{proof} environment of 46 lines). No mathematics is written, repaired or re-derived: the statement and the proof are the article's own lines, in the article's own notation, and no retained line is substituted or re-flowed.  Retained uncounted as the source-local context the proof needs: lines 204--206.  Nothing was omitted from any retained span, so the package carries no omission record.  No retained line contains a citation, so the wrapper carries no bibliography and no bibliography entry.  No retained line contains a figure environment, an image-inclusion command or drawing code, so no image is delivered and the package declares no asset dependency.  Editorial formatting only, in addition: an AMS \texttt{amsart} preamble that restates the article's own declarations and macros used by the retained lines, namely the environments \texttt{proposition} on separate counters, together with the standard packages those lines need.  The retained environments print in this document as Proposition 1 in order of appearance, whereas the article prints the same items with the numbering of its own full document.  The complete native source of the article as distributed in the arXiv e-print for this exact version is bundled unchanged as \texttt{sources/103956/source0.tex}.
\begin{equation}\label{eq333xx}
\varphi_\beta(\alpha) = \bigl( \Tr(\beta\alpha),\; \Tr(\bar{\beta}\alpha) \bigr).
\end{equation}

\begin{proposition}
	The map $\varphi_\beta$ in \eqref{eq333xx} is an $\F_q$-linear isomorphism.
\end{proposition}

\begin{proof}
	For any $k_1, k_2 \in \F_q$ and $\alpha_1, \alpha_2 \in \F_{q^2}$, we have
	\[
	\begin{aligned}
		\varphi_\beta(k_1\alpha_1 + k_2\alpha_2)
		&= \bigl( \Tr(\beta(k_1\alpha_1 + k_2\alpha_2)),\; \Tr(\bar{\beta}(k_1\alpha_1 + k_2\alpha_2)) \bigr) \\
		&= \bigl( k_1\Tr(\beta\alpha_1) + k_2\Tr(\beta\alpha_2),\; k_1\Tr(\bar{\beta}\alpha_1) + k_2\Tr(\bar{\beta}\alpha_2) \bigr) \\
		&= k_1\varphi_\beta(\alpha_1) + k_2\varphi_\beta(\alpha_2),
	\end{aligned}
	\]
	so $\varphi_\beta$ is $\F_q$-linear.
	
	To prove injectivity, assume $\varphi_\beta(\alpha_1) = \varphi_\beta(\alpha_2)$. Then
	\[
	\Tr(\beta(\alpha_1 - \alpha_2)) = 0, \quad\Tr(\bar{\beta}(\alpha_1-\alpha_2)) = 0,
	\]
	which is equivalent to the system
	\begin{equation}\label{eq1}
		\begin{cases}
		\beta(\alpha_1 - \alpha_2) + \bar{\beta}(\alpha_1 - \alpha_2)^q = 0\\
		\bar{\beta}(\alpha_1 - \alpha_2) + \beta(\alpha_1 - \alpha_2)^q = 0
	\end{cases}.
	\end{equation}
	If $\charac(\F_{q^2}) = 2$, then \eqref{eq1} gives
	\[
	\beta(\alpha_1 + \alpha_2) = \bar{\beta}(\alpha_1 + \alpha_2)^q,\quad 
	\bar{\beta}(\alpha_1 + \alpha_2) = \beta(\alpha_1 + \alpha_2)^q,
	\]
	so both $\beta(\alpha_1 + \alpha_2)$ and $\bar{\beta}(\alpha_1 + \alpha_2)$ lie in $\F_q$. Since $\beta$ is primitive, this forces $\alpha_1 = \alpha_2$.
	
	Now suppose $\charac(\F_{q^2}) \neq 2$. Adding and subtracting the two equations in \eqref{eq1} yields
	\[
	\begin{cases}
		(\beta + \bar{\beta})\bigl((\alpha_1 - \alpha_2) + (\alpha_1 - \alpha_2)^q\bigr) = 0\\
		(\beta - \bar{\beta})\bigl((\alpha_1 - \alpha_2) - (\alpha_1 - \alpha_2)^q\bigr) = 0
	\end{cases}.
	\]
	Note that $\beta + \bar{\beta} \neq 0$ and $\beta - \bar{\beta} \neq 0$; otherwise $\beta^{2q-2}=1$ or $\beta^{q-1}=1$, contradicting the fact that $\beta$ has order $q^2-1$ with $q-1 < 2q-2 < q^2-1$. Hence
	\[
	\begin{cases}
		(\alpha_1 - \alpha_2) + (\alpha_1 - \alpha_2)^q = 0\\
		(\alpha_1 - \alpha_2) - (\alpha_1 - \alpha_2)^q = 0
	\end{cases},
	\]
	which implies $2(\alpha_1 - \alpha_2) = 0$ and therefore $\alpha_1 = \alpha_2$ (since $p \neq 2$). Thus $\varphi_\beta$ is injective, and because both spaces have dimension $2$ over $\F_q$, it is an isomorphism.
\end{proof}

\end{document}
